\section{Verified Execution Adapters and Grounded Property Invention}
\label{sec:verified-adapter-property-invention}

The previous executable-falsification system rejected all three original bugs
but could certify only one of three repaired repositories. The principal failure
was not falsification intent; it was unreliable loading and fixture
construction. This stage separated those capabilities into independently
verified execution-adapter acquisition and property invention.

\subsection{Source-attested adapter acquisition}

Given a public issue, original source file and public in-tree tests, the adapter
learner generated a minimal program binding \texttt{AION\_TARGET} to a real
repository symbol. Independent instrumentation recovered the bound object's
source location and required
\begin{equation}
 \operatorname{resolve}(\operatorname{source}(AION\_TARGET))
 = \operatorname{resolve}(p_{\mathrm{declared}}).
\end{equation}
A fabricated marker or substitute implementation could not satisfy this gate.

Marshmallow and pydicom passed immediately. Pvlib required two trace-driven
revisions to avoid unrelated package initialization and load
\texttt{pvlib.iam} with its minimum public dependency. The acquisition session
used five attempts; a deterministic verification replay reconfirmed the three
retained adapters.

\begin{table}[ht]
\centering
\caption{Verified execution-adapter acquisition.}
\label{tab:verified-adapters}
\begin{tabular}{lr}
\hline
Measure & Result \\
\hline
Repositories & 3 \\
Verified adapters & 3/3 \\
Weakest-repository success & 100\% \\
Exact source-path attestation & 100\% \\
Initial attempts & 5 \\
Trace-driven revisions & 2 \\
Unsafe executions / timeouts / live writes & 0 / 0 / 0 \\
Restart retention & 100\% \\
\hline
\end{tabular}
\end{table}

CAU promoted
\texttt{procedure\_verified\_execution\_adapter\_acquisition\_4b5bdc4a2a7e}.

\subsection{Adapter-grounded property invention}

Functional and adversarial properties were then generated inside each attested
adapter. A repository passed only when
\begin{equation}
 \neg T_f(x_{\mathrm{original}}) \land
 T_f(x_{\mathrm{candidate}}) \land
 T_a(x_{\mathrm{candidate}}) \land
 S(T_f,T_a,x_{\mathrm{candidate}}),
\end{equation}
where $T_f$ is the invented functional falsifier, $T_a$ is the invented
adversarial suite and $S$ denotes the static safety gate.

\begin{table}[ht]
\centering
\caption{Adapter-grounded property invention.}
\label{tab:adapter-grounded-properties}
\begin{tabular}{lr}
\hline
Measure & Result \\
\hline
End-to-end repository success & 3/3 (100\%) \\
Weakest-repository success & 100\% \\
Original buggy checkouts rejected & 3/3 \\
Repaired checkouts accepted & 3/3 \\
Adversarial suites accepted & 3/3 \\
Property attempts & 6 \\
Unsafe executions / timeouts / live writes & 0 / 0 / 0 \\
Restart relearning & 0 \\
\hline
\end{tabular}
\end{table}

Criticism removed an over-specific Marshmallow container-shape assumption,
constructed the correct lazy/raw pydicom materialisation path, and eliminated a
pvlib property that contradicted the corrected behavior. It did not expose the
historical human patch or hidden verifier tests. CAU promoted
\texttt{procedure\_adapter\_grounded\_property\_invention\_60fe24a6e28a}.

\subsection{Claim boundary}

The result demonstrates learned execution grounding recomposed with executable
counterexample invention on three public Python repositories. It does not yet
establish repeated-seed stability, three-language transfer, explicit malicious
candidate rejection, conditional memory routing at scale or independently
administered hidden evaluation. Those requirements remain closed gates.
